% GENERATED FILE. Edit canonical lessons, then run npm run generate:books.
% canonical-source-hash: fnv1a64:38492973c13e8b95
% canonical-lessons: FR-C531-bricolage, FR-C531-bricoler, FR-C531-vis, FR-C531-tournevis, FR-C531-percer, FR-R531-revision-530-531, FR-R531-revision-final-531

\chapter{Le bricolage}
\label{ch:fr-le-bricolage}
\hlchaptermodality{531}

\begin{chapteropening}
\textbf{By the end of this chapter:} \emph{I can talk about small repair jobs at home.}

Complete the last lesson of chapter 531: I can talk about small repair jobs at home.
\end{chapteropening}

\section[le bricolage]{le bricolage — DIY, odd jobs}

\label{lesson:FR-C531-bricolage}



\begin{quote}
Before the new one: say the French for to wipe, then the French for tea towel.
\end{quote}

\subsection*{You'll want to know: le bricolage}
\textbf{le bricolage} — \textquotedblleft{}DIY, odd jobs\textquotedblright{}.

\begin{quote}
Le week-end, mon oncle fait du bricolage.
\end{quote}

\begin{quote}
\emph{At the weekend my uncle does DIY.}
\end{quote}

\emph{Un magasin de bricolage} — a DIY shop.

\begin{cousinweb}
\emph{le bricolage} is built on \emph{bricoler}, from Old French \emph{bricole} (EtymDB), once a catapult and later a small job. The earlier origin is not settled.
\end{cousinweb}

\begin{grammarlens}[title={What you've built}]
One more word for telling, explaining and describing at B1.
\end{grammarlens}

\subsection*{Your turn}
\begin{itemize}
\raggedright
  \item \emph{Say it:} \emph{le bricolage}
  \item \emph{Say it:} \emph{le bricolage}, once more
  \item \emph{Say it:} a sentence of your own with \emph{le bricolage}
  \item \emph{Recall:} say the French for to wipe, then the French for tea towel, then say \emph{le bricolage} again
\end{itemize}

\begin{tcolorbox}[breakable,colback=teal!4,colframe=teal!35!black,title={Before you move on}]
What does le bricolage mean? (\textquotedblleft{}DIY, odd jobs\textquotedblright{}.) Say it once more.
\end{tcolorbox}
\section[bricoler]{bricoler — to do DIY, to tinker}

\label{lesson:FR-C531-bricoler}



\begin{quote}
Before the new one: say the French for tea towel, then the French for DIY.
\end{quote}

\subsection*{You'll want to know: bricoler}
\textbf{bricoler} — \textquotedblleft{}to do DIY, to tinker\textquotedblright{}.

\begin{quote}
Ma tante aime bricoler dans son garage.
\end{quote}

\begin{quote}
\emph{My aunt likes tinkering in her garage.}
\end{quote}

\emph{Un bricoleur}, \emph{une bricoleuse} — a handy person.

\begin{cousinweb}
\emph{bricoler} is built on \emph{bricole}, once a catapult and later a small job. The earlier origin is not settled.
\end{cousinweb}

\begin{grammarlens}[title={What you've built}]
One more word for telling, explaining and describing at B1.
\end{grammarlens}

\subsection*{Your turn}
\begin{itemize}
\raggedright
  \item \emph{Say it:} \emph{bricoler}
  \item \emph{Say it:} \emph{bricoler}, once more
  \item \emph{Say it:} a sentence of your own with \emph{bricoler}
  \item \emph{Recall:} say the French for tea towel, then the French for DIY, then say \emph{bricoler} again
\end{itemize}

\begin{tcolorbox}[breakable,colback=teal!4,colframe=teal!35!black,title={Before you move on}]
What does bricoler mean? (\textquotedblleft{}To do DIY, to tinker\textquotedblright{}.) Say it once more.
\end{tcolorbox}
\section[la vis]{la vis — screw}

\label{lesson:FR-C531-vis}



\begin{quote}
Before the new one: say the French for DIY, then the French for to do DIY.
\end{quote}

\subsection*{You'll want to know: la vis}
\textbf{la vis} — \textquotedblleft{}screw\textquotedblright{}.

\begin{quote}
Il manque une vis à cette chaise.
\end{quote}

\begin{quote}
\emph{This chair is missing a screw.}
\end{quote}

The \emph{s} is sounded: \emph{vis} rhymes with \emph{fils}.

\begin{cousinweb}
\emph{la vis} is inherited from Latin \emph{vītis} \textquotedblleft{}a vine\textquotedblright{} (EtymDB; Lewis \& Short): the screw is probably named for the spiral of the vine's tendrils.
\end{cousinweb}

\begin{grammarlens}[title={What you've built}]
One more word for telling, explaining and describing at B1.
\end{grammarlens}

\subsection*{Your turn}
\begin{itemize}
\raggedright
  \item \emph{Say it:} \emph{la vis}
  \item \emph{Say it:} \emph{la vis}, once more
  \item \emph{Say it:} a sentence of your own with \emph{la vis}
  \item \emph{Recall:} say the French for DIY, then the French for to do DIY, then say \emph{la vis} again
\end{itemize}

\begin{tcolorbox}[breakable,colback=teal!4,colframe=teal!35!black,title={Before you move on}]
What does la vis mean? (\textquotedblleft{}Screw\textquotedblright{}.) Say it once more.
\end{tcolorbox}
\section[le tournevis]{le tournevis — screwdriver}

\label{lesson:FR-C531-tournevis}



\begin{quote}
Before the new one: say the French for to do DIY, then the French for screw.
\end{quote}

\subsection*{You'll want to know: le tournevis}
\textbf{le tournevis} — \textquotedblleft{}screwdriver\textquotedblright{}.

\begin{quote}
Pour ouvrir la radio, il a pris un petit tournevis.
\end{quote}

\begin{quote}
\emph{To open the radio, he took a small screwdriver.}
\end{quote}

Literally turn-screw. The \emph{s} is sounded.

\begin{cousinweb}
\emph{le tournevis} is built from \emph{tourner} \textquotedblleft{}to turn\textquotedblright{}, from Latin \emph{tornāre}, and \emph{vis} \textquotedblleft{}screw\textquotedblright{}, inherited from Latin \emph{vītis} \textquotedblleft{}a vine\textquotedblright{} (EtymDB).
\end{cousinweb}

\begin{grammarlens}[title={What you've built}]
One more word for telling, explaining and describing at B1.
\end{grammarlens}

\subsection*{Your turn}
\begin{itemize}
\raggedright
  \item \emph{Say it:} \emph{le tournevis}
  \item \emph{Say it:} \emph{le tournevis}, once more
  \item \emph{Say it:} a sentence of your own with \emph{le tournevis}
  \item \emph{Recall:} say the French for to do DIY, then the French for screw, then say \emph{le tournevis} again
\end{itemize}

\begin{tcolorbox}[breakable,colback=teal!4,colframe=teal!35!black,title={Before you move on}]
What does le tournevis mean? (\textquotedblleft{}Screwdriver\textquotedblright{}.) Say it once more.
\end{tcolorbox}
\section[percer]{percer — to drill, to pierce}

\label{lesson:FR-C531-percer}



\begin{quote}
Before the new one: say the French for screw, then the French for screwdriver.
\end{quote}

\subsection*{You'll want to know: percer}
\textbf{percer} — \textquotedblleft{}to drill, to pierce\textquotedblright{}.

\begin{quote}
Mon voisin perce un trou dans le mur.
\end{quote}

\begin{quote}
\emph{My neighbour is drilling a hole in the wall.}
\end{quote}

Spelling: \emph{nous perçons}. \emph{Une perceuse} — a drill.

\begin{cousinweb}
\emph{percer} is inherited from Old French \emph{percier} (EtymDB), from Vulgar Latin \emph{pertūsiāre}, built on Latin \emph{pertundere} \textquotedblleft{}to bore through\textquotedblright{} (Lewis \& Short).
\end{cousinweb}

\begin{grammarlens}[title={What you've built}]
One more word for telling, explaining and describing at B1.
\end{grammarlens}

\subsection*{Your turn}
\begin{itemize}
\raggedright
  \item \emph{Say it:} \emph{percer}
  \item \emph{Say it:} \emph{percer}, once more
  \item \emph{Say it:} a sentence of your own with \emph{percer}
  \item \emph{Recall:} say the French for screw, then the French for screwdriver, then say \emph{percer} again
\end{itemize}

\begin{tcolorbox}[breakable,colback=teal!4,colframe=teal!35!black,title={Before you move on}]
What does percer mean? (\textquotedblleft{}To drill, to pierce\textquotedblright{}.) Say it once more.
\end{tcolorbox}
\section[(dialogue)]{Loisirs et maison}

\label{lesson:FR-R531-revision-530-531}



\begin{quote}
Say the French for dishes and (kitchen) sink.
\end{quote}

\subsection*{Your turn}
Say these aloud:

\begin{itemize}
\raggedright
  \item dishes, (kitchen) sink, to rinse, to wipe, tea towel
  \item DIY, to do DIY, screw, screwdriver, to drill
\end{itemize}

\begin{tcolorbox}[breakable,colback=teal!4,colframe=teal!35!black,title={Before you move on}]
Dishes? (\textbf{la vaisselle}.) (kitchen) sink? (\textbf{l'évier}.) To rinse? (\textbf{rincer}.) To wipe? (\textbf{essuyer}.) Tea towel? (\textbf{le torchon}.) DIY? (\textbf{le bricolage}.) To do DIY? (\textbf{bricoler}.) Screw? (\textbf{la vis}.) Screwdriver? (\textbf{le tournevis}.) To drill? (\textbf{percer}.)
\end{tcolorbox}
\section[(dialogue)]{Le bricolage — once more}

\label{lesson:FR-R531-revision-final-531}



\begin{quote}
Say the French for DIY and to do DIY.
\end{quote}

\subsection*{Your turn}
\begin{itemize}
\raggedright
  \item \emph{Say it:} DIY, to do DIY, screw, screwdriver, to drill
\end{itemize}

\begin{tcolorbox}[breakable,colback=teal!4,colframe=teal!35!black,title={Before you move on}]
DIY? (\textbf{le bricolage}.) To do DIY? (\textbf{bricoler}.) Screw? (\textbf{la vis}.) Screwdriver? (\textbf{le tournevis}.) To drill? (\textbf{percer}.)
\end{tcolorbox}
